\section{A denominator bound without balanced multiplicities}
\label{sec:denominator}

The contour numerator is normalized by a vacuum polynomial. Its lower
bound must hold even when nearly every site has the same height. We
prove that uniform statement before analyzing the numerator. The proof
uses a positive Pfaffian integral and a Cauchy interpolation formula;
there is no comparison of unequal multiplicities to a balanced list.


Define
\[
 H_*(u)=
 \frac{e^{2u}+2de^u+1}{(e^u-1)e^{u/2}}
 \tanh(\gamma u/2),
\]
with its removable value at \(u=0\).

\begin{theorem}[Uniform denominator lower bound]
\label{thm:06-denominator}
For every even integer \(N\ge2\), let all \(N\) site logs, including the marks,
be purely imaginary with heights in \(\{0,\beta\}\).  Uniformly in
their two multiplicities, as \(N\to\infty\),
\begin{equation}
 \log\left|
 \frac{\widetilde h_N^2}
 {(\gamma/2)^N\prod_{i<j}H_*(a_i-a_j)^2}
 \right|
 \ge \left(\frac5{24}-o(1)\right)\log N .
 \label{eq:06-denominator}
\end{equation}
\end{theorem}

\begin{proof}
A common shift leaves the ratio unchanged.  Center the sites so
that \(|\Im a_i|\le\beta/2\), and initially take distinct imaginary
sites anywhere in this interval.  The Pfaffian formula for \(h_N\)
and the reference Pfaffian product formula show that, before
dividing by \((\gamma/2)^N\), the cleared ratio is
\[
 \left(\frac{\operatorname{Pf}[f(a_i-a_j)]}
 {\operatorname{Pf}[f_0(a_i-a_j)]}\right)^2,\qquad
 f(u)=\frac{\sinh u}{\cosh u+d},\quad
 f_0(u)=\tanh(\gamma u/2).
\]
The odd kernels have the convergent integral representations
\[
\begin{split}
 f_0(u)&=\int_0^\infty
 (e^{iku}-e^{-iku})\,
 \frac{dk}{i\gamma\sinh(3\beta k)},\\
 f(u)&=\frac{\gamma}{2}\int_0^\infty
 (e^{iku}-e^{-iku})
 \left(1+\frac1{2\cosh(2\beta k)}\right)
 \frac{dk}{i\gamma\sinh(3\beta k)} .
\end{split}
\]
They hold on all differences being used.  They follow, for example,
by Fourier transforming the integrable derivatives and then
integrating; oddness fixes the integration constant.
At \(k=0\) the difference of exponentials cancels the measure's
simple pole, and at infinity the available strip widths give
absolute convergence.

Put \(m=N/2\).
Expanding the \(m\)-th exterior power of these kernels gives a
Pfaffian integral over \(0<k_1<\cdots<k_m\), whose integrand is the
determinant with column pairs
\((e^{ik_j a_i},e^{-ik_j a_i})\), times the product of the
corresponding measures.  This is the exterior-power form of the determinant--Pfaffian
integration identity \cite{deBruijn1955}. Taking a difference within each column
pair proves absolute convergence near zero.
Writing \(a_i=iy_i\), all entries of this determinant are real
exponentials.  With \(y_i\)'s ordered, the determinant never
vanishes when the \(2m\) exponents are distinct: a nontrivial linear
combination of \(r\) exponentials with distinct real exponents has
at most \(r-1\) real zeros, by dividing by its first exponential
and applying Rolle's theorem inductively.  After fixing the column
order, the determinant therefore has a fixed sign on the chamber.
The sign argument is a real-exponential instance of the total-positivity
principle; compare the variation-diminishing background in
\cite{Schoenberg1950PFII}. Here the elementary Rolle induction gives
the strict sign needed for these specific nodes. Removing its constant phase gives a probability measure
\(\mathbb P_0\) from the reference integrand.

Define
\[
 S(k)=\log\left(1+\frac1{2\cosh(2\beta k)}\right).
\]
The normalized Pfaffian ratio is
\(\mathbb E_0\exp(\sum_{j=1}^m S(k_j))\).
Consequently Jensen's inequality gives
\begin{equation}
 \log\left|
 \frac{\widetilde h_N^2}
 {(\gamma/2)^N\prod_{i<j}H_*(a_i-a_j)^2}
 \right|
 \ge 2\mathbb E_0\sum_{j=1}^m S(k_j).
 \label{eq:06-jensen}
\end{equation}
To evaluate the right side, deform the reference spectral measure
by \(1+tS(k)\), and denote its odd kernel by \(f_t\).
The logarithmic derivative at zero of
\((\operatorname{Pf}f_t/\operatorname{Pf}f_0)^2\) is exactly this
right side.

Adding any constant all-ones matrix to an even-size skew matrix
does not change its determinant.  Indeed the rank-one determinant
identity gives the multiplier
\(1+c\,\mathbf1^{\mathsf t}A^{-1}\mathbf1=1\) for invertible skew
\(A\), and the identity follows in general by polynomial
continuation.  Add the constant \(1+tS(0)\) to \(f_t\), and
conjugate the difference kernel by \(e^{-\gamma u/2}\).
At \(t=0\) this gives
\[
 K(u)=\operatorname{sech}(\gamma u/2),\qquad
 \widehat K(k)=\frac{2\pi}{\gamma}\operatorname{sech}(3\beta k).
\]
The derivative kernel \(J\) has transform
\[
 \widehat J(k)=\widehat K(k)T(k),\qquad
 T(k)=S(k-i\gamma/2).
\]
For justification, the derivative before conjugation is the
convolution of \(f_0+1\) with the inverse transform of \(S\):
the added constant is exactly its constant term \(S(0)\).
The function \(1+1/(2\cosh(2\beta k))\) is holomorphic and nonzero in
\(|\operatorname{Im} k|<\pi/(4\beta)\): the first zero of
\(2\cosh(2\beta k)\) is on that strip boundary, and the first zero
of \(2\cosh(2\beta k)+1\) has imaginary height \(\pi/(3\beta)\).
Thus \(S\) extends to this strip as the holomorphic logarithm that
is real on the real axis. It decays exponentially
on smaller closed strips, and $\gamma/2=\pi/(6\beta)$ lies inside
this strip. Thus the conjugation shifts the multiplier as displayed.

Let \(K[a]\) be the matrix \(K(a_i-a_j)\), and put
\(L_0=\gamma^{-1}\log N\).
We will show, uniformly in the node positions,
\begin{equation}
 \frac{\operatorname{Tr}(K[a]^{-1}J[a])}{2L_0}
 \longrightarrow\int_{\R}T(k)\,\frac{dk}{2\pi}.
 \label{eq:06-trace-limit}
\end{equation}
This uses interpolation directly and hence requires no symmetry of
the two site multiplicities.
The Cauchy determinant formula gives
\[
\begin{split}
 G(z,w)&:=K[z,a]K[a]^{-1}K[a,w]\\
 &=K(z-w)\bigl(1-B(z)B(w)\bigr),\\
 B(z)&=\prod_{i=1}^N
 \frac{e^{\gamma z}-e^{\gamma a_i}}
 {e^{\gamma z}+e^{\gamma a_i}} .
\end{split}
\]
This identity follows by applying the bordered Cauchy determinant
formula to the interpolation error; its one-point factors are those
in the displayed product.

Choose \(\beta/2<h<\beta\), and define a smooth integrable function
\(g\) by
\[
 \widehat g(-k)
 =\frac{T(k)}{\widehat K(k)\,2\cosh(2hk)} .
\]
The symbol and all its derivatives decay exponentially because
\(2h>\beta\).  Expanding the finite interpolation sum and applying
Fourier inversion in \(x\) gives
\begin{equation}
 \operatorname{Tr}(K[a]^{-1}J[a])
 =\sum_{\pm}\int_{\R}\int_{\R}
 G(x+r/2\pm ih,x-r/2\mp ih)\,g(r)\,dr\,dx .
 \label{eq:06-interpolation-trace}
\end{equation}
In fact the \(x\)-integration of the product of its two \(K\)'s
has symbol \(\widehat K(k)^2e^{\mp2hk}\);
summing signs and integrating \(r\) leaves
\(\widehat K(k)T(k)\), which is precisely \(J\).
All shifts stay within the regular strips.

For $z=x\pm ih$ and $a_i=iy_i$, $|y_i|\le\beta/2$, put
$c_0=\cos(\gamma(h+\beta/2))>0$. If $r=e^{\gamma x}$ and
$\vartheta_i=\gamma(\pm h-y_i)$, the squared modulus of the
$i$th factor of $B$ is
\[
 \frac{r^2+1-2r\cos\vartheta_i}{r^2+1+2r\cos\vartheta_i},
 \qquad \cos\vartheta_i\ge c_0.
\]
Every factor has modulus at most one. The inequality
$r/(r+1)^2\ge\min(r,r^{-1})/4$ and telescoping products give,
uniformly in the nodes,
\[
 |B(z)|\le\exp(-c_1N e^{-\gamma|x|}),\qquad
 |1-B(z)|\le C_1N e^{-\gamma|x|}
 \quad\hbox{in the two tails}.
\]
The second bound at the negative end uses that \(N\) is even.
It follows that for each fixed \(r\)
\[
 \frac1{2L_0}\int_{\R}
 [1-B(x+r/2\pm ih)B(x-r/2\mp ih)]\,dx\longrightarrow1.
\]
Indeed in \(|x|\le(1-\delta)L_0\) the product tends uniformly to
zero, while outside \(|x|\le(1+\delta)L_0\) the tail integral is
negligible.  The transition region costs \(O(\delta L_0)\).
These estimates are uniform in the node positions.
The normalized absolute integral is also bounded independently of
\(r\): use
\(|1-B(z)B(w)|\le |1-B(z)|+|1-B(w)|\) and translate each bound.
Since \(g(r)K(r\pm2ih)\) is integrable, dominated convergence in
\eqref{eq:06-interpolation-trace} proves
\eqref{eq:06-trace-limit}.

The same uniform estimates pass to coincident nodes.
One can take divided differences in the finite interpolation
matrices, or first take limits in the Pfaffian inequality:
after clearing the Vandermonde factors, the ratios are continuous
by polynomiality. Explicitly, on an allowed difference $u=iy$,
$|y|\le\beta$,
\[ H_*(iy)=(\cos y+d)\tan(\gamma y/2)/\sin(y/2), \]
with removable value $\gamma(1+d)>0$ at zero; it is finite and
nonzero on the entire interval. We pass to the limit in this cleared
scalar ratio, not in the inverse interpolation matrix. The uniform
positive lower bound obtained
above prevents loss to zero in this limit.

Finally shift the integral of \(T\) back to the real line.
For \(k>0\), setting \(y=e^{-2\beta k}\) gives
\[
 S(k)=\log(1-y^3)-\log(1-y)-\log(1+y^2).
\]
Integration of the absolutely integrable logarithmic series yields
\[
 \int_{\R}S(k)\,dk=\frac{5\pi^2}{72\beta}.
\]
Combining \eqref{eq:06-jensen} and \eqref{eq:06-trace-limit}, and
using \(\gamma\beta=\pi/3\), gives
\[
 2L_0\,\frac{5\pi^2}{72\beta}\,\frac1{2\pi}
 =\frac5{24}\log N,
\]
which proves \eqref{eq:06-denominator}.
\end{proof}

